%%
%% This is a placeholder manuscript shell for WWW'27 (Research Track — Security and Privacy).
%% Refresh this manuscript before final submission by syncing with docs/paper/manuscript.tex
%% from the research repository.
%%
\documentclass[sigconf,anonymous,review]{acmart}

\acmConference[WWW '27]{The ACM Web Conference 2027}{May 10--14, 2027}{Dublin, Ireland}
\acmYear{2027}
\acmDOI{}
\acmISBN{}

\title[What a Client Fetches]{What a Client Fetches: Default-Path Exposure of Models Published on the Web}

\author{Anonymous Author(s)}
\affiliation{%
  \institution{Anonymous Institution}
  \country{Anonymous Country}
}
\email{anonymous@example.org}

\begin{abstract}
Machine learning models are increasingly published on the Web through public hubs and fetched by client applications on default execution paths. This paper investigates a fundamental provenance question for the Security and Privacy track of The Web Conference: when a derived model declares an upstream parent on a public hub, does the client fetch an artefact published by that parent, or does the derivation introduce unvetted execution hazards? We formally distinguish file presence from default-path exposure, where default client loaders fetch hazardous weight formats or dynamic configurations without explicit overrides. Evaluating a complete depth-one lineage census ($N = 48{,}247$ declared edges across 16 foundational base models) on the Hugging Face Hub, 0 of the 16 parents exhibit default-path serialisation exposure, establishing that inheritance of default-path exposure is not identified. Across child repositories, only 3.17\% (95\% CI: [3.01\%, 3.33\%]) exhibit default-path serialisation exposure, whereas 54.90\% maintain legacy pickle files bypassed by default loaders. Downstream mutations outnumber repairs by 8.5x, and 11.22\% of declared children launder upstream licences. Model supply chain risks on the Web arise primarily from downstream modification rather than parent inheritance.
\end{abstract}

\begin{CCSXML}
<ccs2012>
   <concept>
       <concept_id>10002978.10003022.10003023</concept_id>
       <concept_desc>Security and privacy~Software security engineering</concept_desc>
       <concept_significance>500</concept_significance>
       </concept>
   <concept>
       <concept_id>10002951.10003260.10003282</concept_id>
       <concept_desc>Information systems~Web applications</concept_desc>
       <concept_significance>300</concept_significance>
       </concept>
 </ccs2012>
\end{CCSXML}

\ccsdesc[500]{Security and privacy~Software security engineering}
\ccsdesc[300]{Information systems~Web applications}

\keywords{Web security, provenance, model hub, default-path exposure, serialisation, supply chain}

\begin{document}

\maketitle

\section{Introduction}
Public model repositories such as the Hugging Face Hub have become the primary Web infrastructure for distributing pre-trained machine learning checkpoints. Applications routinely download and execute remote model weights over HTTPS using automated library interfaces such as \texttt{AutoModel.from\_pretrained()}. When a client fetches a model that reports an upstream parent, does the client fetch an artefact derived safely from that parent, or does the derivation introduce severe execution hazards along the Web supply chain?

We address this question by evaluating the byte streams served across the Web fetch boundary. We distinguish between static file presence and default-path exposure: modern loaders prioritise safe formats (e.g., safetensors), rendering legacy pickle files inactive unless explicit overrides are specified.

\section{Web Architecture and Default-Path Exposure}
We define the threat model and trust boundaries governing model distribution on the Web:
\begin{enumerate}
    \item \textbf{Arbitrary Deserialisation Execution:} Python pickle bytecode embedded in legacy checkpoints can trigger arbitrary code execution upon deserialisation.
    \item \textbf{Remote Code Execution via Configuration:} Custom code bindings declared in \texttt{auto\_map} configurations execute when clients load models with remote code execution enabled.
    \item \textbf{Licence Laundering:} Downstream fine-tunes stripping upstream intellectual property constraints or commercial restrictions.
\end{enumerate}

\section{Census Methodology}
We conduct an exhaustive census of $N = 48{,}247$ parent--child lineage edges across 16 major foundational base model families on the Hugging Face Hub. Our automated pipeline evaluates file listings, serialisation formats, configuration fields, and licence compatibility.

\section{Empirical Findings}
Key results from our census include:
\begin{itemize}
    \item \textbf{Parent Exposure:} 0 of the 16 base models exhibit default-path serialisation exposure; inheritance of default-path serialisation exposure is not identified.
    \item \textbf{Child Exposure vs Presence:} Only 3.17\% of child models ($1{,}528$ of $48{,}247$) exhibit default-path serialisation exposure, whereas 54.90\% ($26{,}486$) exhibit pickle file presence.
    \item \textbf{Mutation Dynamics:} Downstream developers introduce pickle files into safe lineages on $21{,}486$ edges (52.77\%), outnumbering repairs ($2{,}529$ edges) by an introduction-to-repair ratio of 8.5:1.
    \item \textbf{Licence Laundering:} 11.22\% of declared children ($3{,}019$ edges) declare incompatible licences, while 44.21\% ($21{,}328$ edges) drop licence metadata entirely.
\end{itemize}

\section{Implications for Web Security}
These findings demonstrate that risks in the Web model ecosystem stem primarily from downstream repository mutations rather than upstream parent inheritance. Client-side verification gates and package provenance must be enforced at the Web fetch boundary.

\section{Conclusion}
This study formalises default-path exposure for Web-distributed models, demonstrating that safe defaults mitigate legacy file presence while highlighting downstream mutation as the primary attack surface.

\bibliographystyle{ACM-Reference-Format}
\bibliography{references}

\end{document}
